% !TEX root = ../VLDB-2026-DAGSmith-Dependency-Aware-Rewriting-for-dbt-Style-SQL-Pipelines.tex
\section{Related Work}
\label{sec:related}

The most closely related lines of work are database query optimization, source-to-source SQL rewriting, multi-query optimization, materialized-view recommendation, workflow scheduling, and data lineage or clone analysis for analytics pipelines. \tool differs from these lines by treating the explicit SQL pipeline DAG itself as the optimization object, allowing it to restructure node boundaries and reason jointly about semantics, materialization, and refresh frequency.

Database query optimizers remain essential but orthogonal: they choose physical plans for one submitted SQL statement. \tool operates before that point, changing the maintained SQL pipeline program that produces those statements. Source-to-source SQL rewriters such as \emph{SlabCity}~\cite{slabcity}, \emph{GenRewrite}~\cite{genrewrite}, \emph{WeTune}~\cite{wetune}, and \emph{R-Bot}~\cite{rbot} also operate before the database optimizer, but they optimize one statement at a time and do not change the dependency graph around it.

Multi-query optimization and common-subexpression sharing reason across multiple statements, but typically over flat workloads and mostly through syntactic overlap. They do not usually exploit output contracts, declared refresh cadences, materialization choices, or stable developer-maintained node boundaries. Materialized-view selection optimizes persistence decisions, but does not usually rewrite the graph to move computation across freshness boundaries or remove intermediate contracts. Workflow systems such as Airflow~\cite{airflowDags}, Dagster~\cite{dagsterAssets}, Prefect~\cite{prefectTasks}, Argo~\cite{argoDag}, Luigi~\cite{luigiDocs}, SQLMesh~\cite{sqlmeshOverview}, and Databricks Workflows~\cite{databricksJobs} expose dependencies and schedules, but their optimizers primarily decide ordering, invalidation, or execution policy rather than semantic SQL rewrites. \tool targets the gap between these areas: dependency-aware source-to-source optimization for recurring SQL pipeline DAGs.
